\section{Ejercicio I}
\subsection{Banco Ahorro Famsa}
En Banco Ahorro Famsa se observa un crecimiento acelerado de la cartera de crédito entre 2016 y mayo de 2020. La cartera total bruta pasó de 14,180 millones de pesos en 2016 a 27,443 millones de pesos en mayo de 2020. Sin embargo, este crecimiento estuvo acompañado por un deterioro significativo en la calidad de los activos, ya que la cartera vencida aumentó de 1,719 millones a 4,951 millones de pesos en el mismo periodo. 

El índice de morosidad también mostró un deterioro importante. En 2016, el índice de morosidad fue de 12.12\%, mientras que para mayo de 2020 aumentó a 18.04\%. Esto indica que una proporción cada vez mayor de los créditos otorgados por el banco no estaba siendo recuperada en tiempo. La presión fue especialmente visible en los créditos al consumo, particularmente en los créditos personales, los cuales concentraban la mayor parte de la cartera vencida. Asimismo, la estimación preventiva para riesgos crediticios aumentó de 1,898 millones de pesos en 2016 a 6,029 millones de pesos en mayo de 2020. Aunque la cobertura de cartera vencida se mantuvo por encima del 100\%, este incremento no debe interpretarse como una señal de fortaleza, sino como evidencia de un mayor riesgo crediticio.
En este caso, la cobertura pasó de 110.41\% en 2016 a 121.77\% en mayo de 2020. Esto significa que el banco estaba constituyendo reservas superiores al monto de la cartera vencida; sin embargo, el aumento constante de dichas reservas refleja que la institución necesitaba reconocer pérdidas esperadas cada vez mayores.Por su parte, el capital contable creció entre 2016 y 2019, al pasar de 3,298 millones a 5,516 millones de pesos. No obstante, para mayo de 2020 disminuyó a 5,121 millones de pesos. Esta caída coincide con un resultado neto negativo de -567 millones de pesos, lo que sugiere que el deterioro de la cartera y el aumento de las reservas comenzaron a erosionar directamente la posición patrimonial del banco.En conjunto, la información muestra que Banco Ahorro Famsa no presentó únicamente un problema de crecimiento de cartera, sino un deterioro progresivo en la calidad de sus activos. El banco aumentó su cartera de crédito, pero al mismo tiempo acumuló una proporción creciente de cartera vencida. Esto obligó a incrementar las estimaciones preventivas y redujo la rentabilidad, hasta afectar el capital contable. Por ello, la revocación de Banco Ahorro Famsa puede explicarse como resultado de una combinación de alta morosidad, presión sobre reservas, pérdidas netas y debilitamiento patrimonial.

\subsection{Banco Bicentenario}

En Banco Bicentenario se observa un deterioro financiero acelerado entre diciembre de 2013 y junio de 2014. A diciembre de 2013, la institución registraba una cartera total bruta de 852.9 millones de pesos, integrada por 768.5 millones de cartera vigente y 84.4 millones de cartera vencida. Para junio de 2014, la cartera total disminuyó a 767.0 millones de pesos; sin embargo, la cartera vencida aumentó de manera significativa hasta 377.9 millones de pesos. Esto muestra que, aunque el tamaño total de la cartera se redujo, la proporción de créditos problemáticos creció de forma muy marcada.


En diciembre de 2013, el índice de morosidad fue de 9.90\%. Para junio de 2014, este indicador aumentó a 49.27\%. Esto significa que casi la mitad de la cartera de crédito del banco se encontraba vencida. Este deterioro es especialmente grave porque indica una pérdida acelerada de calidad de los activos crediticios y una fuerte incapacidad de recuperación de los créditos otorgados.

La estimación preventiva para riesgos crediticios también aumentó de forma considerable. En diciembre de 2013 fue de 60.3 millones de pesos, mientras que para junio de 2014 se ubicó aproximadamente en 359.0 millones de pesos. Este crecimiento refleja que el banco tuvo que reconocer mayores pérdidas esperadas por el deterioro de su cartera.


En diciembre de 2013, la cobertura fue de 71.45\%, lo cual indica que las reservas preventivas no alcanzaban a cubrir completamente la cartera vencida. Para junio de 2014, la cobertura aumentó a aproximadamente 95.00\%. Aunque este incremento muestra un mayor reconocimiento del riesgo crediticio, no fue suficiente para corregir el deterioro estructural del banco, ya que la cartera vencida había crecido de forma desproporcionada.

El capital contable también muestra una caída importante. En diciembre de 2013, Banco Bicentenario tenía un capital contable de 281.2 millones de pesos. Para junio de 2014, este se redujo a 110.0 millones de pesos. Esta disminución refleja que las pérdidas y el deterioro de la cartera comenzaron a erosionar directamente la base patrimonial de la institución.

La situación patrimonial se agravó aún más al observar el índice de capitalización. En junio de 2014, Banco Bicentenario tenía un índice de capitalización de 9.99\%, ubicado en categoría II. Sin embargo, para julio de 2014, el índice de capitalización cayó a 2.98\%, colocándolo en categoría V, la categoría más grave de alertas tempranas. Esto significa que el banco quedó muy por debajo de los niveles mínimos de capital requeridos para operar de forma segura.

En conjunto, Banco Bicentenario muestra un caso de deterioro financiero muy rápido. La institución no solo enfrentó un aumento fuerte de la cartera vencida, sino también una reducción de su cartera total, un crecimiento acelerado de las estimaciones preventivas, una caída del capital contable y una pérdida severa de capitalización. La causa principal de su quiebra puede explicarse por la combinación de mala calidad de cartera, insuficiencia de reservas en la etapa inicial del deterioro, pérdidas patrimoniales y una caída crítica del índice de capitalización. Por ello, Banco Bicentenario representa un caso claro de insolvencia derivada del deterioro crediticio y de la incapacidad del banco para absorber las pérdidas generadas por su cartera.

\subsection{Consultoría Internacional Banco}

En el caso de CIBanco, la evolución financiera muestra un comportamiento distinto al de Banco Ahorro Famsa y Banco Bicentenario. Entre 2016 y 2024 no se observa un deterioro crediticio severo ni una caída sostenida del capital contable. Por el contrario, la institución mantuvo una cartera relativamente estable, niveles moderados de cartera vencida o deteriorada y un crecimiento gradual de su capital.

Entre 2016 y 2021, la cartera total bruta pasó de 9,068 millones de pesos a 13,579 millones de pesos, mientras que la cartera vencida aumentó de 111 millones a 361 millones. Aunque hubo crecimiento en términos absolutos, el índice de morosidad se mantuvo bajo, pasando de 1.22\% en 2016 a 2.66\% en 2021. Esto indica que CIBanco no presentaba una crisis crediticia tradicional durante ese periodo.

A partir de 2022, bajo el formato IFRS9, la cartera se clasifica por etapas de riesgo. Para mantener la comparación, se considera la cartera en etapa 1 y etapa 2 como cartera vigente o no deteriorada, mientras que la etapa 3 se utiliza como aproximación de cartera deteriorada. Con esta metodología, el índice de deterioro fue de 2.96\% en 2022, 3.23\% en 2023 y 3.16\% en 2024. Estos niveles muestran que, hasta 2024, el deterioro crediticio seguía siendo moderado.

El capital contable también mostró fortaleza relativa, al pasar de 2,061 millones de pesos en 2016 a 5,702 millones en 2024. Asimismo, la estimación preventiva para riesgos crediticios aumentó de 162 millones en 2016 a 532 millones en 2024, manteniendo una cobertura razonable frente a la cartera deteriorada.

El cambio más importante aparece en agosto de 2025. La cartera total bruta cayó de 15,046 millones de pesos en 2024 a 6,941 millones, mientras que la cartera deteriorada aumentó de 476 millones a 563 millones. Como resultado, el índice de deterioro subió a 8.11\%. Además, la cobertura bajó de 111.76\% en 2024 a 64.65\% en agosto de 2025, lo que indica que las reservas dejaron de cubrir completamente la cartera deteriorada.

A pesar de este deterioro en 2025, la situación de CIBanco no parece corresponder a una quiebra clásica por morosidad acumulada. Hasta 2024, sus indicadores eran relativamente estables y su capital contable seguía creciendo. Por ello, la salida de CIBanco del sistema bancario debe entenderse principalmente como resultado de un choque regulatorio y reputacional. En 2025, FinCEN señaló a CIBanco como institución de preocupación principal en materia de lavado de dinero, lo que afectó su viabilidad operativa y la confianza en la institución.

En conjunto, CIBanco representa un caso distinto al de bancos quebrados por deterioro interno de cartera. Aunque en 2025 se observa una contracción fuerte de la cartera y una menor cobertura de reservas, el detonante principal fue de naturaleza regulatoria, reputacional y operativa, más que una crisis crediticia tradicional.
\subsection{Intercam Banco}
En el caso de Intercam Banco, la evolución financiera muestra una trayectoria relativamente sólida entre 2016 y 2024, pero con un cambio abrupto en 2025. A diferencia de Banco Ahorro Famsa y Banco Bicentenario, Intercam no presenta durante la mayor parte del periodo una crisis clásica por deterioro sostenido de cartera. Por el contrario, hasta 2024 se observa crecimiento de la cartera total, aumento del capital contable y niveles de morosidad o deterioro generalmente moderados.

Entre 2016 y 2021, la cartera total bruta pasó de 6,363 millones de pesos a 11,835 millones de pesos. Durante ese mismo periodo, la cartera vencida tuvo un comportamiento variable: aumentó de 147 millones en 2016 a 551 millones en 2018, pero posteriormente disminuyó hasta 199 millones en 2021. Esto permitió que el índice de morosidad bajara de 5.91\% en 2018 a 1.68\% en 2021, lo que indica una mejora en la calidad de la cartera antes de la transición al formato IFRS9.

A partir de 2022, bajo el formato IFRS9, la cartera se clasifica por etapas de riesgo. Para mantener la comparación con los años anteriores, se considera la cartera en etapa 1 y etapa 2 como cartera vigente o no deteriorada, mientras que la cartera en etapa 3 se utiliza como aproximación de cartera deteriorada. Bajo esta metodología, Intercam registró una cartera total bruta de 16,396 millones de pesos en 2022, 19,676 millones en 2023 y 25,147 millones en 2024. Esto muestra una expansión importante de la cartera antes del deterioro observado en 2025.

El índice de deterioro también se mantuvo en niveles relativamente contenidos entre 2022 y 2024. En 2022 fue de 3.13\%, en 2023 disminuyó a 1.92\% y en 2024 aumentó a 3.48\%. Aunque este último dato muestra un incremento en la cartera deteriorada, no representa por sí solo una crisis crediticia comparable con la de Banco Ahorro Famsa o Banco Bicentenario.

La estimación preventiva para riesgos crediticios también creció conforme aumentó la cartera. En 2016 fue de 176 millones de pesos y para 2021 llegó a 697 millones. Posteriormente, bajo IFRS9, se ubicó en 1,023 millones en 2022, 912 millones en 2023 y 1,348 millones en 2024. La cobertura de cartera vencida o deteriorada fue variable, pero en general se mantuvo en niveles suficientes en los últimos años previos a la crisis, alcanzando 199.03\% en 2022, 241.27\% en 2023 y 154.06\% en 2024.

El capital contable muestra una trayectoria creciente hasta 2024. Pasó de 1,101 millones de pesos en 2016 a 4,061 millones en 2021. Posteriormente aumentó a 6,013 millones en 2022, 8,531 millones en 2023 y 11,052 millones en 2024. Esto sugiere que, antes del choque de 2025, Intercam mantenía una posición patrimonial amplia y no mostraba una erosión progresiva de capital como la observada en otros bancos quebrados por deterioro financiero interno.

El cambio más importante ocurre en 2025. Para el corte reportado, la cartera total bruta cae de 25,147 millones de pesos en 2024 a 891 millones de pesos. Al mismo tiempo, la cartera deteriorada se mantiene en 657 millones, lo que eleva el índice de deterioro a 73.74\%. Este cambio no parece reflejar únicamente un aumento ordinario de morosidad, sino una contracción extrema de la cartera total, probablemente asociada a la pérdida de viabilidad operativa, reducción de operaciones o salida de activos tras el choque regulatorio y reputacional.

En ese mismo corte, el capital contable cae de 11,052 millones de pesos en 2024 a 4,023 millones de pesos. Aunque la estimación preventiva de 891 millones de pesos cubre la cartera deteriorada en 135.62\%, la magnitud de la contracción de cartera y la caída del capital muestran un cambio estructural en la institución. Por ello, el deterioro de 2025 debe interpretarse con cautela: no parece ser el resultado de una acumulación gradual de cartera vencida, sino de un evento disruptivo que alteró de forma abrupta la operación del banco.

En conjunto, Intercam Banco representa un caso similar al de CIBanco. Hasta 2024, sus indicadores financieros no muestran una quiebra clásica por morosidad acumulada; por el contrario, se observa crecimiento de cartera, aumento del capital contable y cobertura razonable de reservas. Sin embargo, en 2025 se presenta una contracción fuerte de la operación y un deterioro relativo muy elevado. Por esta razón, la salida o crisis de Intercam debe explicarse principalmente por factores regulatorios, reputacionales y operativos, más que por un deterioro crediticio tradicional acumulado.